%%%%%%%%%%%%%%%%%%%%%%%%%%%%%%%%%%%%%%%%%%%%%%%%%%%%%%%%%%%%%%%%%%%%%%%%
\chapter{Introduzione}
%%%%%%%%%%%%%%%%%%%%%%%%%%%%%%%%%%%%%%%%%%%%%%%%%%%%%%%%%%%%%%%%%%%%%%%%

\begin{definition}
  Denotiamo l'insieme dei numeri naturali da $ 1 $ a $ n \in \N $ come $ [n] = \left\{ 1,\dots,n
  \right\} $, e il prodotto cartesiano tra $ [n] $ e $ [m] $ come $ [n] \otimes [m] =
  \left\{ (i,j) : i \in [n], j \in [m] \right\} $.
\end{definition}

\begin{definition}
  Un \emph{tensore} $ \Tn{X}  $ di ordine $ d $ è un array multidimensionale con entrate definite sul campo $
  \K $, che denotiamo come $  \Tn{X} \in \Kmsiz{\K} $. Indicizziamo gli elementi di un tensore come
  $ \Tn{X}(\miwc) $, dove $ (\miwc) \in \mdom$.
\end{definition}


\todo{ADD: Esempi di tensori}



\begin{definition}
  Date due matrici $ A \in \R^{m_1 \times m_2} $ e $ B \in \R^{m_2 \times n_2} $, il loro prodotto
  di Kroneker è la matrice 
  \[
  C = A \krn B = 
  \begin{bmatrix} 
    a_{11}B & a_{12}B & \ldots & a_{1n}B \\ 
    a_{21}B & a_{22}B & \ldots & a_{2n}B \\ 
    \vdots & \vdots & \ddots & \vdots \\ 
    a_{m1}B & a_{m2}B & \ldots & a_{mn}B \\ 
  \end{bmatrix}  \in  \R^{m_1m_2 \times n_1n_2}, \text{dove } c_{kl} = a_{i_1,j_1}b_{i_2j_2},
  .\] 
  dove $ k = (i_1-1)m_2+i_2 $ e $ l = (j_1-1)n_2 + j_2 $
  \todo{riscrivere con L*}
  
\end{definition}

\begin{definition}[Tensor slice]
  Dato $ \Tn{x} \in \Rmsiz $ chiamiamo slice (o fetta) di $ \Tn{x} $ ogni matrice ottenuta da $
  \Tn{x} $ fissando tutti gli indici tranne due. Le slices ottenute fissando gli ultimi $ d-2 $ indici sono dette frontal slices.
\end{definition}

\begin{definition}[Ordinamento naturale]
  L'ordinamento naturale (natural ordering) su $ \mdom! $ è la mappa biunivoca
\begin{equation*}
  \begin{aligned}
    \L \colon \mdom! &\to [\prod_{j=1}^d n_j]\\
    (\miwc) &\mapsto 1 + s_1 (i_1-1) + \ldots + s_d(i_d-1) = 1 + \sum_{j=1}^d s_j(i_j - 1),
  \end{aligned}
\end{equation*}
dove $ s_j = n_1 \cdots n_{j-1} = \prod_{h=1}^{j-1}n_h $ è detto passo $ j- $esimo.
L'inversa di $ \L $ è data dalla mappa
\begin{equation*}
  \begin{aligned}
    \L^{-1} \colon [\prod_{j=1}^d n_j] &\to \mdom!\\
    l &\mapsto \left(1 + \left\lfloor \frac{ (l-1)\mod(n_1s_1)}{s_1}\right\rfloor, \ldots, 1 +
    \left\lfloor \frac{ (l-1)\mod(n_ds_d)}{s_d} \right\rfloor \right)
  \end{aligned}
\end{equation*}
\end{definition}
\begin{remark}
  La mappa $ \L $ ha come inversa $ \L^{-1} $, sono anche chiamate linearizzazione e tensorizzazione
  degli indici, rispettivamente, in quanto permettono di passare da uni indice tensoriale a uno
  lineare, e viceversa.
\end{remark}
\begin{proof}
Segue direttamente dalla definizione di $ \L $ e dalle proprietà del modulo, infatti,
\[
\begin{aligned}
  \L^{-1}(\L(\miwc)) &= \L^{-1}(1 + s_1(i_1-1)) + \ldots + s_d(i_d-1)) \\
                     &= \left( 1+ \left\lfloor \frac{\sum_{i=1}^d s_j(i_j-1) \mod(n_1s_1)}{s_1}
                     \right\rfloor, \ldots,   1+ \left\lfloor \frac{\sum_{i=1}^d s_j(i_j-1)
                   \mod(n_ds_d)}{s_d} \right\rfloor \right)\\
                     &= \left( 1+ \left\lfloor \sum_{i=1}^d \frac{s_j}{s_1}(i_j-1) \mod(n_1s_1)
                     \right\rfloor, \ldots,   1+ \left\lfloor \sum_{i=1}^d \frac{s_j}{s_d}(i_j-1)
                   \mod(n_ds_d) \right\rfloor \right)\\
                     &= \left( \miwc \right),
\end{aligned}
\]
dove nell'ultimo passaggio abbiamo usato che $ s_j = \prod_{k=1}^d n_k$ e il fatto che $
\left\lfloor x \right\rfloor = 0 $ per $ x \in [0,1) $.
Viceversa,
\[
\begin{aligned}
  \L(\L^{-1}(l)) &= \L\left(1 + \left\lfloor \frac{ (l-1)\mod(n_1s_1)}{s_1}\right\rfloor, \ldots, 1 +
    \left\lfloor \frac{ (l-1)\mod(n_ds_d)}{s_d} \right\rfloor\right) = \\
                 &=1 + s_1\left(\left\lfloor \frac{(l-1) \mod(n_1s_1)}{s_1} \right\rfloor \right) + \ldots +
                 s_d\left(\left\lfloor \frac{(l-1) \mod(n_ds_d)}{s_d} \right\rfloor\right)
\end{aligned}
\] 
\todo{missing end}
\end{proof}

Analogamente alle matrici, i tensori vengono salvati in memoria come vettori, risulta quindi
naturale definire la seguente oprazione
\begin{definition}
  Dato un tensore $ \Tn{X}$, possiamo definire l'operazione che trasforma un tensore in vettore,
  tramite la mappa di vettorizzazione 
\begin{equation*}
\begin{aligned}
  \vec \colon \Rmsiz &\to \R^{\prod_{j=1}^d n_j}\\
  \Tn{x} &\mapsto x\\
  \tikz[baseline=1ex, scale=0.4]{
    % Front face
    \draw (0,0) rectangle (1,1);
    % Top face
    \draw (0,1) -- (0.3,1.3) -- (1.3,1.3) -- (1,1);
    % Right face
    \draw (1.3,1.3) -- (1.3,0.3) -- (1,0);
  } 
  &\mapsto 
  \tikz[baseline=-1ex, scale=0.4]{
    % Tall rectangle
    \draw (0,-1.5) rectangle (0.3, 1.5);
  }
\end{aligned}
\end{equation*}
 dove $ x_l = \Tn{X}(\L^{-1}(l;\miwc[n])) $.
\end{definition}

\begin{proposition}
  Siano $ \Mx{A} \in \K^{m \times n} $, $ \Mx{B} \in \K^{m \times n}$ ed $ \Mx{X} \in \K^{n \times m}$. Allora
  \begin{equation*}
    \vec({\Mx{A}\Mx{X}\Mx{B}}) = (B^{\top} \krn A)\vec(X)
  \end{equation*}
\end{proposition}
\begin{proof}
  Scrivendo $ B = [b_1, \ldots, b_n] $ ed $ X = [x_1, \ldots, x_n] $, la $ k- $esima colonna di $
  AXB $ è 
  \begin{equation*}
    \begin{aligned}
      (AXB)_{:,k} &= AXb_k = A \sum_{i=1}^m x_i b_{i,k} = \left[ b_{1,k}A ,\ldots, b_{n,k}A \right]
      \begin{bmatrix} x_1 \\ \vdots \\ x_m \end{bmatrix} = ([b_{1,k}, \ldots, b_{m,k}] \krn A)
      \vec(X).
    \end{aligned}
  \end{equation*}
  Impilando insieme le colonne otteniamo che
  \begin{equation*}
  \vec(AXB) = \begin{bmatrix} (AXB)_{:,1} \\(AXB)_{:,2}\\ \vdots \\ (AXB)_{:,n} \end{bmatrix} =
  \begin{bmatrix} b_1^{\top} \krn A \\ b_2^{\top} \krn A \\ \vdots \\ b_n^{\top} \krn A
  \end{bmatrix} \vec(X) = (B^{\top} \krn A)\vec(X).
  \end{equation*}
\end{proof}

\begin{proposition}
  Siano $ \Mx{A} \in \R^{m \times p} $ e $ \Mx{B} \in \R^{n \times q} $. Allora, 
  \[
  \Mx{A} \krn \Mx{B} = \Mx{P}(\Mx{B} \krn \Mx{A})\Mx{Q}^{\top},
  .\] 
  dove $ \Mx{P} $ e $ \Mx{Q} $ sono matrici di permutazione $ (m,n)- $ e $ (p,q)- $ perfect shuffle, rispettivamente.
\end{proposition}
\begin{proof}
\todo{Potremmo riadattare: THE VEC-PERMUTATION MATRIX, THE VEC OPERATOR AND KRONECKER PRODUCTS: A REVIEW}

\end{proof}

\section{La fattorizzazione CP}
\todo{Short introduzione che leghi caso matriciale con tensoriale.}
\begin{definition}[Canonical Polyadic Decomposition (CP)]
  Diciamo che un tensore $ \Tn{x} \in \Rmsiz $ ammette una decomposizione CP di rango $ r\in \N $ se
  esistono delle matrici $ U_1, \ldots, U_d $ tali che $ U_h = \left[ u_1^{(h)}, \ldots, u_r^{(h)}\right]\in R^{n_h \times r} $ tali per cui
  valgono le seguenti tre condizioni equivalenti:
  \begin{enumerate}[label=(\roman*)]
    \item $ \Tn{x} = \sum_{j=1}^r u_j^{(1)} \topp \ldots \topp u_j^{(d)} = \sum_{i=1}^r
      \bigtopp_{h=1}^d u_j^{(h)}$.
    \item $ \vec(\Tn{x}) = \sum_{j=1}^r u_j^{(d)} \krn \ldots \krn u_j^{(1)} = \sum_{i=1}^r
      \bigkrn_{h=d}^1 u_j^{(h)}$.
    \item $\Tn{x}(\miwc) = \sum_{j=1}^r U_1(i_1,j)U_2(i_2,j)\ldots U_d(i_d,j) = \sum_{j=1}^r
      \prod_{h=1}^d U_h(i_h,j)$.
  \end{enumerate}
  In tal caso, scriviamo $ \Tn{x} = \dsqr{U_1, \ldots, U_d} $.
\end{definition}

\subsection{Formulare il prodotto di matrici con la CP}
